\documentclass[11pt]{article}
\usepackage[T1]{fontenc}
\usepackage{lmodern}
\usepackage[letterpaper,margin=1.02in]{geometry}
\usepackage{amsmath,amssymb,amsthm}
\usepackage[hidelinks]{hyperref}
\usepackage{url}
\usepackage{etoolbox}

\providecommand{\paperauthor}{}
\newtheorem{theorem}{Theorem}
\newtheorem{lemma}[theorem]{Lemma}
\newcommand{\id}[1]{\texttt{\detokenize{#1}}}
\newcommand{\A}{A}
\newcommand{\singlerow}[1]{\put(0,#1){\line(1,0){1}}}
\newcommand{\crossrows}[1]{\put(0,#1){\line(1,1){1}}\put(1,#1){\line(-1,1){1}}}
\newcommand{\stripfour}[1]{%
  \setlength{\unitlength}{0.205in}%
  \begin{picture}(1.32,3.30)(-.16,-.15)
    \linethickness{.65pt}#1%
    \multiput(0,0)(0,1){4}{\circle*{.17}}%
    \multiput(1,0)(0,1){4}{\circle*{.17}}%
  \end{picture}}
\hypersetup{pdftitle={Perfect matchings with two axial edges on the king's grid}}
\setlength{\emergencystretch}{1.2em}
\raggedbottom

\begin{document}
\begin{center}
  {\LARGE\bfseries Perfect matchings with two axial edges\\[2pt]
    on the king's grid\par}
  \ifdefempty{\paperauthor}{}{\medskip{\large\paperauthor\par}}
\end{center}

\begin{abstract}
On a rectangular king's grid, an axial edge joins two cells sharing a
side. We prove that the number of perfect matchings with exactly two axial
edges on a \(2m\times2n\) board is \(mn(m+n+2)/2\). On the square this
gives \(n^2(n+1)\), as conjectured by Ry{\l}ko. A row and column parity
argument forces the two edges to be parallel. The free boundary then
splits the board into independent strips, exactly one of which contains
both edges. A Lean~4 formalization verifies that the strip construction
enumerates each matching exactly once.
\end{abstract}

\section{The two-edge coefficient}

In the king's graph on a rectangular board, two distinct cells are
adjacent when their row and column coordinates each differ by at most
one. All edges lie within the board. A perfect matching pairs every cell
once, and two diagonal edges may cross. Call horizontal and vertical edges
\emph{axial}. Let \(\A_j(m,n)\) count the perfect matchings of a board
with \(2m\) rows and \(2n\) columns containing exactly \(j\) axial edges.

\begin{theorem}\label{thm:main}
For all nonnegative integers \(m,n\),
\[
  \A_2(m,n)
   =n\binom{m+1}{2}+m\binom{n+1}{2}
   =\frac{mn(m+n+2)}{2}.
\]
The first summand counts matchings with two horizontal edges and the
second counts those with two vertical edges.
\end{theorem}

Ry{\l}ko's array \(T(n,k)\) counts matchings of the \(2n\times2n\)
square with \(2k\) diagonal edges. He recorded the last two entries of each
row as conjectures, supported by computations through \(n=10\)
\cite{rylko}. Since a matching on the square has \(2n^2\) edges, the
penultimate entry is \(T(n,n^2-1)=\A_2(n,n)\) for \(n\geq1\).
Theorem~\ref{thm:main} gives the conjectured value \(n^2(n+1)\).

\section{From a rectangle to one strip}

\begin{lemma}\label{lem:parity}
The horizontal edges of a perfect matching occur equally often in odd
and even rows. The vertical edges occur equally often in odd and even
columns.
\end{lemma}

\begin{proof}
Each nonhorizontal edge covers one cell in an odd row and one in an even
row. The two classes have the same size, so the horizontal edges cover
equally many cells in each class. Each such edge covers two cells in a
single row. Transposing the board proves the assertion for columns.
\end{proof}

\newpage
\begin{proof}[Proof of Theorem~\ref{thm:main}]
An empty board has no two-edge configuration, so assume \(m,n\geq1\).
By Lemma~\ref{lem:parity} both axial edges are horizontal or both are
vertical. Suppose they are horizontal. Every cell in the first column must
take a partner in the second. These \(2m\) distinct partners exhaust
that column. Removing both columns and repeating partitions the matching
into \(n\) independent strips of width two.

In each strip consider the two top cells. If they are matched to one
another, delete that single row (write \(S\)). Otherwise each has only one
remaining neighbor, the diagonally opposite cell in the next row.
The two forced edges form a cross. Delete those two rows (write \(X\)).
Iterating gives a word in \(S,X\) of total height \(2m\), where \(S\)
has height one and \(X\) has height two. Conversely any such word
fills the strip from the top down. Each \(S\) contributes one horizontal
edge and each \(X\) contributes none.

Every strip has even height, so its number of letters \(S\) is even.
Exactly one strip therefore contains both occurrences of \(S\).
The other strips each have the unique word \(X^m\). In the exceptional
strip the two identical letters \(S\) and the \(m-1\) identical letters
\(X\) occupy \(m+1\) positions, giving \(\binom{m+1}{2}\) words.
Choosing the strip contributes the factor \(n\). Transposing gives
\(m\binom{n+1}{2}\) matchings in the vertical case. These two classes are
disjoint, and their sum is the claimed count.
\end{proof}

\begin{center}
  \begin{tabular}{ccc}
   \stripfour{\singlerow{3}\singlerow{2}\crossrows{0}} &
   \stripfour{\singlerow{3}\crossrows{1}\singlerow{0}} &
   \stripfour{\crossrows{2}\singlerow{1}\singlerow{0}}\\[2pt]
   \(SSX\) & \(SXS\) & \(XSS\)
  \end{tabular}\\[4pt]
  \small The three exceptional strips of height four, read from the top.
\end{center}

For example, on a \(4\times6\) board the count is \(3\cdot3+2\cdot6=21\).
The single-strip enumeration also has a classical permutation description.
If \(\pi(i)\) is the right-column partner row of the left cell in row
\(i\), then \(|\pi(i)-i|\leq1\). The letters \(S\) and \(X\) are the fixed
points and adjacent transpositions of a Fibonacci permutation.
Diaconis, Graham and Holmes count permutations of this kind on
\(N\) letters with \(k\) transpositions by \(\binom{N-k}{k}\)
\cite[\S2.4, Propositions 2.1 and 2.4]{dgh}.
Taking \(N=2m\) and \(k=m-1\) recovers the exceptional-strip factor.
The parity and boundary arguments determine how that factor enters the
count for the whole board.

The same strip description gives \(\A_0(m,n)=1\), the other end entry of
Ry{\l}ko's array on the square. This uniqueness also follows from earlier
tiling work. Diagonal edges preserve checkerboard color. On either color,
rotating and rescaling the centers turns corner adjacency into ordinary
side adjacency on a mixed Aztec rectangle. The checkerboard construction
on the square is given by Ciucu \cite[\S3]{ciucu}. For the rectangle,
start with the white cells of a \((2m+1)\times(2n+1)\) board with black
corners and remove the white cells on two adjacent boundary sides. This
leaves the white cells of the even rectangle and gives Lee's mixed Aztec
rectangle, up to a rigid motion. Lee records its unique domino tiling
\cite[\S6, Figure~14]{lee}. Reflection exchanges the two colors, giving
one matching on each color and thus one on the board.

\newpage
\section{Formal verification}

The accompanying Lean~4 project defines a legal matching by a list of
king edges whose endpoints cover each board cell exactly once. It then
associates to that list its undirected edge relation \(R\). Reordering
the list or reversing any of its edges leaves \(R\) unchanged. The
predicate \id{IsTwoTiling} additionally requires the sum of the horizontal
and vertical edge counts to be two.

In the namespace \id{KingEnumeration}, the theorem
\id{exactly_two_rectangle} produces a list \(L\) of relations with no
duplicates, proves that \(R\in L\) exactly when \(R\) is such a tiling,
and proves the natural-number identity
\[
  2\,\mathrm{length}(L)=n\,m(m+1)+m\,n(n+1).
\]
The proof builds edges from strip words, recovers the strip encoding of
any legal matching in the horizontal case, and shows that decoding
recovers the original undirected edge relation. Transposition supplies
the other orientation. The corollary \id{exactly_two_square} gives the
length directly as \(n^2(n+1)\).

The project uses Lean 4.19.0 and its bundled libraries. From the project
root, \id{lake build} builds the library and \id{python scripts/verify.py}
rebuilds the sources in a fresh directory and reports the axioms used by
the named theorems. The rectangular theorem uses only the standard Lean
axioms \id{propext}, \id{Classical.choice} and \id{Quot.sound}. An
independent finite check, \id{python scripts/check_small_boards.py}, pairs
the first uncovered cell with each possible free neighbor and checks all
25 rectangles with \(0\leq m,n\leq4\). The accompanying source archive
includes the Lean project and both verification scripts under
\id{anc/king-grid/}.

\paragraph{Use of generative AI.}
OpenAI Codex was used to search for related work, develop the counting
argument, implement the Lean proofs and reproduction scripts, and draft
the manuscript.

\begin{thebibliography}{9}
\bibitem{ciucu}
M. Ciucu, \href{https://doi.org/10.1007/s10801-007-0099-7}{Symmetry
classes of spanning trees of Aztec diamonds and perfect matchings of odd
squares with a unit hole}, \emph{J. Algebraic Combin.} \textbf{27}
(2008), 493--538.

\bibitem{dgh}
P. Diaconis, R. Graham and S.~P. Holmes,
\href{https://doi.org/10.1214/lnms/1215090070}{Statistical problems
involving permutations with restricted positions}, in
\emph{State of the Art in Probability and Statistics},
IMS Lecture Notes--Monograph Series \textbf{36} (2001), 195--222.

\bibitem{lee}
Y.-L. Lee, \href{https://doi.org/10.37236/10913}{An extension of the
Lindstr\"om--Gessel--Viennot theorem}, \emph{Electron. J. Combin.}
\textbf{29} (2022), no.~2, Paper No.~2.41.

\bibitem{rylko}
D. Ry{\l}ko, \href{https://oeis.org/A401673}{Entry A401673},
\emph{The On-Line Encyclopedia of Integer Sequences} (2026).
\end{thebibliography}
\end{document}
